%=========== ΛΥΣΗ - 20ο ΘΕΜΑ Γ ===========
\begin{thema}{Γ}
\begin{erwthma}
\item Το πεδίο ορισμού της $f(x) = \sqrt{x+1}$ είναι το $[-1, +\infty)$. Στο $x_0 = 0$, η συνάρτηση είναι παραγωγίσιμη (αφού $0 > -1$). 
Υπολογίζουμε την τιμή της συνάρτησης στο $0$:
\[ f(0) = \sqrt{0+1} = 1 \]
Άρα το σημείο επαφής είναι το $A(0, 1)$.
Η παράγωγος της $f$ για $x > -1$ είναι:
\[ f'(x) = (\sqrt{x+1})' = \frac{1}{2\sqrt{x+1}} \cdot (x+1)' = \frac{1}{2\sqrt{x+1}} \]
Η κλίση της εφαπτομένης ευθείας $\varepsilon$ στο $x_0 = 0$ είναι:
\[ \lambda = f'(0) = \frac{1}{2\sqrt{0+1}} = \frac{1}{2} \]
Η εξίσωση της εφαπτομένης είναι:
\[ y - f(0) = f'(0)(x - 0) \implies y - 1 = \frac{1}{2}x \implies y = \frac{1}{2}x + 1 \]

\item Για τη μελέτη της κυρτότητας, υπολογίζουμε τη δεύτερη παράγωγο για $x > -1$:
\begin{align*}
f''(x) &= \left( \frac{1}{2\sqrt{x+1}} \right)' = \frac{1}{2} \left( (x+1)^{-\frac{1}{2}} \right)' \\
&= \frac{1}{2} \left( -\frac{1}{2} \right) (x+1)^{-\frac{3}{2}} \cdot (x+1)' \\
&= -\frac{1}{4(x+1)^{\frac{3}{2}}} = -\frac{1}{4\sqrt{(x+1)^3}}
\end{align*}
Επειδή για κάθε $x > -1$ ισχύει $(x+1)^3 > 0$, ο παρονομαστής είναι θετικός, οπότε:
\[ f''(x) < 0 \text{ για κάθε } x \in (-1, +\infty) \]
Άρα η συνάρτηση $f$ είναι \textbf{κοίλη} (στρέφει τα κοίλα κάτω) στο $[-1, +\infty)$.
Εφόσον η $f$ είναι κοίλη, η γραφική της παράσταση βρίσκεται \textbf{«κάτω»} από οποιαδήποτε εφαπτομένη της (με εξαίρεση το σημείο επαφής).
Χρησιμοποιώντας την εφαπτομένη $\varepsilon$ στο $x_0 = 0$ που βρήκαμε στο (α), ισχύει για κάθε $x \ge -1$:
\[ f(x) \le y_{\varepsilon} \implies \sqrt{x+1} \le \frac{1}{2}x + 1 \]

\item Στο ερώτημα (β) αποδείξαμε ότι:
\[ \sqrt{x+1} \le \frac{x}{2} + 1 \text{ για κάθε } x \ge -1 \]
Η ανισότητα αυτή ισχύει προφανώς και στο κλειστό διάστημα $[0, 3]$. 
Επειδή οι συναρτήσεις είναι συνεχείς και δεν ταυτίζονται στο $[0,3]$ (η ισότητα ισχύει μόνο στο $0$), ολοκληρώνοντας και τα δύο μέλη προκύπτει:
\[ \int_{0}^{3} \sqrt{x+1} \d x < \int_{0}^{3} \left( \frac{x}{2} + 1 \right) \d x \]
(η ανισότητα ισχύει και με «$\le$» όπως ζητάει το ερώτημα, προφανώς).
Υπολογίζουμε το ολοκλήρωμα του δεξιού μέλους:
\begin{align*}
\int_{0}^{3} \left( \frac{x}{2} + 1 \right) \d x &= \left[ \frac{x^2}{4} + x \right]_{0}^{3} = \left( \frac{3^2}{4} + 3 \right) - (0 + 0) \\&= \frac{9}{4} + 3 = \frac{9}{4} + \frac{12}{4} = \frac{21}{4}
\end{align*}
Κατά συνέπεια, αποδείχθηκε ότι:
\[ \int_{0}^{3} \sqrt{x+1} \d x \le \frac{21}{4} \]

\item Ζητείται ο υπολογισμός του ορίου:
\[ L = \lim_{x \to 0} \frac{1}{2\sqrt{x+1} - 2 - x} \]
Από την ανισότητα που αποδείξαμε στο Γ2, έχουμε για κάθε $x \ge -1$:
\[ \sqrt{x+1} \le \frac{x}{2} + 1 \]
Πολλαπλασιάζοντας με το $2$ (που είναι θετικό) παίρνουμε:
\[ 2\sqrt{x+1} \le x + 2 \implies 2\sqrt{x+1} - x - 2 \le 0 \]
Γνωρίζουμε ότι η ισότητα ισχύει \textbf{μόνο} στο σημείο επαφής, δηλαδή για $x = 0$.
Επομένως, για κάθε $x$ "κοντά" στο $0$ με $x \neq 0$, ισχύει:
\[ 2\sqrt{x+1} - 2 - x < 0 \]
Καθώς το $x \to 0$, ο παρονομαστής $2\sqrt{x+1} - 2 - x$ τείνει στο $0$ (αφού $2\sqrt{1} - 2 - 0 = 0$). 
Όμως, λόγω της αυστηρής ανισότητας που δείξαμε παραπάνω, ο παρονομαστής προσεγγίζει το μηδέν παραμένοντας διαρκώς \textbf{αρνητικός}. 
Δηλαδή, ο παρονομαστής τείνει στο $0^-$.
Ο αριθμητής είναι θετικός ($1 > 0$), άρα το όριο είναι:
\[ \lim_{x \to 0} \frac{1}{2\sqrt{x+1} - 2 - x} = -\infty \]
\end{erwthma}
\end{thema}
%=========== ΤΕΛΟΣ ΛΥΣΗΣ - 20ο ΘΕΜΑ Γ ===========
